% Optional narrative sketch; not included by paper/_main.tex.
% Updated alongside the formal manuscript to avoid retaining superseded claims.
\section*{Narrative Sketch: Persistence at Cosmic Scale}

\paragraph{The organizing question.}
How might a finite-capacity system maintain persistent distinctions as its internal complexity changes? Papers I--V provide candidate substrate geometry, updates and operational constraints. Applying them to the universe requires a cosmological state and readout, rather than identifying every finite-system analogy with a physical law.

\paragraph{Initialization.}
A proposed chiral selection reduces the active channel labels from 32 to 16. If all input labels must remain distinguishable, reversible bookkeeping retains additional information; a binary record suffices for that label-space example. Inflation, an arrow of time and the primordial spectrum require a dynamical account of that process. The candidate tilt $31/32$ is a channel-complement formula, and zero source alone does not force zero tensor initial data. The two upstream CP formulas give different estimates under the stipulated baryogenesis conversion; neither formula supplies baryon-number-changing dynamics.

\paragraph{Load and cosmological inference.}
The revised budget is $\mathcal L=1+S_Q+S_T+S_M=1491/1376$. Under the proposed response rules it raises the draft early-inferred present-day Hubble rate and lowers the clustering input, with an exact invariant product. An absolute calibration, expansion history and growth equations remain necessary. These are inference/readout hypotheses, not established idle and loaded epochs with independently derived rates.

\paragraph{Saturation.}
The local index law gives a positive subcapacity clock speed and a zero speed limit at saturation. Its nonlinear averaging properties are proved, but an exterior metric and global causal structure are not. Finite reversible storage requires additional outcomes or capacity for fresh independent input. This constrains possible information handling without choosing buffering, reflection, Hawking radiation or an evaporation law. Entropy counting gives $A\log q/a$; the physical normalization must still be derived.

\paragraph{Dark-sector interpretation.}
Mode counts require energy weights and occupancies before becoming density ratios. Total-to-active and inactive-to-active counts are different. A load-dependent density yields a conditional equation of state only after a continuity law is specified. The old common-scale contextual-neutrino spectrum fails the upstream splitting checks, so it cannot be presented as a resolution of the joint particle/cosmology problem.

\paragraph{Purpose.}
The persistence perspective suggests where to seek operational laws and which information must be retained. The next goal is to derive those laws and observables, then test them jointly. Finite capacity alone does not force every proposed cosmological effect or its numerical coefficient.
